\section{Translating Relations in Alloy}


\begin{enumerate}
	\item Sig are treated as relations of arity 1. The representation of a general relation is an extension to arity n, with the type being the same in TLA+
	\item Each relation is associated with a VARIABLE, just like a sig
	\item Relations are defined by their product type, and constraints. 
	\item The product type is used for the membership constraints on the VARIABLE
	\item The implicit constraints emerge from quantifiers used in the field declaration
	\item All fields are treated as global, the translator does not support the declaration of multiple fields with the same name
	\item There are two types of fields: ones which do not refer to any other field, and ones which do
	\item Define Head(f) as the sig inside which the field is declared, and Tail(f) as the list of sigs and fields within the decl
	\item The goal of the computation of the product type is to define a function Type(f) which is a list of sigs, such that when their cross product is taken, we get the product type.
	\item When computing product types, all implicit constraints are ignored
	\item Algorithm for translation: in the base case, where Tail(f) has no sigs, Type(f) is Head(f)::Tail(f)
	\item In the case where Tail(f) starts with a field g, Type(f) is Type(g)++fold(map(Tail(Tail(f), Type),++)
	\item In the case where Tail(f) starts with a sig S, Type(f) is S::fold(map(Tail(f), Type),++)
	\item In cases where the decl elements are set expressions, There is an ambiguity in what happens if Head(f) involves a field
	\item The above are restricted to cases where f and g live in the same sig and share a Head. If the Heads are different, the rules change
\end{enumerate}


